\section{Contemporary Evidence: The Thermodynamic Collapse in Real-Time}
\label{sec:contemporary}

\subsection{Framing: Mechanism Demonstration, Not Prevalence Claim}

This section examines documented AI/ML deployment failures and adjacent automation failures that share the maintenance-extraction dynamic central to this paper's thermodynamic framework. The cases presented are \textit{illustrative mechanism-probes}, selected according to explicit criteria: high-stakes domain deployment, external validation or regulatory finding, litigation with documented outcomes, or executive acknowledgment of reversal. They are not intended as prevalence claims; systematic survey of AI deployment outcomes remains beyond this paper's scope.

Two distinct failure modes emerge from the evidence, both amenable to thermodynamic interpretation:

\begin{enumerate}
\item \textbf{Capability failures}: distribution shift between training and deployment environments, hallucination as architectural feature, tacit knowledge gaps that pattern extraction cannot bridge, maintenance entropy in complex systems.

\item \textbf{Governance failures}: deployment of automation to deny rights or care, opacity as accountability evasion, incentive structures that tolerate high error rates when contestation costs exceed user resources.
\end{enumerate}

These categories interact but are analytically distinct. Governance failures can dominate even when capability is adequate; capability failures can occur even in well-governed deployments. Both, however, share a common thermodynamic signature: extraction of value from systems without corresponding investment in maintenance, validation, or human oversight.

\subsection{The Implementation Crisis: Quantifying System Brittleness}

The transition from pilot to production reveals systematic brittleness across AI deployments. S\&P Global Market Intelligence reports that 42\% of surveyed companies abandoned the majority of their AI initiatives by early 2025, a sharp increase from 17\% the previous year, with organizations scrapping an average of 46\% of proof-of-concepts before reaching production \citep{spglobal2025}. The RAND Corporation's analysis of AI project outcomes found failure rates exceeding 80\%---approximately twice that of conventional IT projects \citep{rand2024}.

These figures find corroboration across industry analyses. \citet{mckinsey2025} reports that fewer than 10\% of deployed use cases progress beyond pilot stage, with only 1\% of companies achieving ``mature'' deployment. Boston Consulting Group's survey of 1,000 executives across 59 countries found that only 26\% developed capabilities to move beyond proof of concept; 74\% remained in cycles of failed experiments \citep{bcg2024}. Gartner forecasts that 30\% of generative AI projects will be abandoned after proof of concept by end of 2025, citing poor data quality, inadequate risk controls, and unclear business value \citep{gartner2024ai}.

The pattern suggests predictable brittleness under conditions of complexity:

\begin{enumerate}
\item Initial enthusiasm underestimates maintenance requirements
\item Pilot success in controlled conditions masks distribution shift
\item Scaling attempts encounter complexity that vectorized training cannot navigate
\item Abandonment as error rates exceed organizational tolerance
\item Repetition with new initiatives facing identical constraints
\end{enumerate}

\subsection{The Deskilling Spiral: The Mechanism Linking Signal to Cases}

The macro-level brittleness documented above operates through a specific mechanism at the human-system interface. \citet{rintakahila2023} documented this dynamic with empirical precision in accounting firms: ``Cognitive automation leads to complacency and reduces mindfulness in tasks, gradually eroding essential skills.'' The degradation follows thermodynamic logic: capabilities not maintained through active investment decay regardless of whether replacement systems function adequately.

The cycle accelerates through documented stages:

\begin{enumerate}
\item Automation reduces human practice frequency (investment withdrawal)
\item Reduced practice erodes skill maintenance (entropy accumulation)
\item Skill erosion increases automation dependency (brittleness amplification)
\item Dependency accelerates further automation pressure (feedback loop)
\item System achieves simultaneous deskilling and automation inadequacy (collapse conditions)
\end{enumerate}

This mechanism explains why the cases that follow exhibit similar failure signatures despite spanning different domains. Whether the context is clinical decision support, customer service, or legal research, the extraction of human cognitive work without corresponding investment in human capability maintenance creates the brittleness that macro statistics capture. The Epic Sepsis Model's 88\% false positive rate trains clinicians to ignore algorithmic warnings; UnitedHealth's nH Predict overrides physician judgment by design; legal AI hallucinations persist because verification skills atrophy when research feels automated.

The deskilling spiral thus provides the causal bridge between implementation-level abandonment rates and the specific failure modes documented in the cases below. Organizations extract cognitive value from human workers, encode partial patterns into automated systems, then discover that neither the degraded human capabilities nor the brittle automated systems can navigate genuine complexity. The thermodynamic interpretation: you cannot extract expertise without investing in its maintenance and expect the resulting system to remain stable.

\subsection{Capability Failures: The Training-Deployment Gap}

\subsubsection{Case: Knight Capital---Automation Entropy Precursor}

\textit{Note: This case predates contemporary AI/ML systems but illustrates the maintenance-extraction dynamic that recurs across automation deployments.}

\textbf{What happened}: On August 1, 2012, Knight Capital Group lost \$440 million in 45 minutes when a technician failed to copy updated code to one of eight servers. The eighth server retained dormant ``Power Peg'' code from 2005---a testing algorithm designed to buy high and sell low. When triggered at market open, this unmaintained code executed over 4 million trades across 154 stocks, accumulating \$3.5 billion in unwanted positions \citep{sec2013knight, nytimes2012knight}.

\textbf{Failure mode}: Maintenance entropy. Dormant code accumulated technical debt that entropy collected catastrophically when the system encountered unexpected activation conditions.

\textbf{Thermodynamic interpretation}: Extraction of operational value from automated systems without corresponding investment in maintenance, testing, and change-control creates brittleness that manifests unpredictably under stress.

\textbf{Ethical stake}: Market integrity; systemic risk from inadequately maintained trading infrastructure; Knight Capital's subsequent acquisition by Getco LLC transferred consequences to shareholders and employees.

\subsubsection{Case: Epic Sepsis Model---Vendor Claims vs. Validated Reality}

\textbf{What happened}: Researchers at Michigan Medicine conducted external validation of Epic's widely-implemented sepsis prediction algorithm, published in \textit{JAMA Internal Medicine} \citep{wong2021sepsis}. The algorithm achieved area under curve (AUC) of 0.63 versus vendor-claimed 0.76--0.83. At Michigan's threshold, sensitivity reached only 33\% (missing 67\% of sepsis cases) with positive predictive value of 12\% (88\% false positive rate). The system flagged 18\% of all hospitalized patients while missing two-thirds of actual sepsis cases.

\textbf{Failure mode}: Distribution shift between training environment and deployment context; opacity of proprietary validation preventing external scrutiny.

\textbf{Thermodynamic interpretation}: Performance metrics extracted from controlled training distributions do not transfer to the thermodynamic complexity of actual clinical environments. As the researchers noted: ``The increase and growth in deployment of proprietary models has led to an underbelly of confidential, non--peer-reviewed model performance documents that may not accurately reflect real-world model performance.''

\textbf{Ethical stake}: Patient safety; alert fatigue degrading clinician responsiveness; opacity preventing informed institutional decisions about deployment.

\subsubsection{Case: IBM Watson for Oncology---The \$4 Billion Capability Ceiling}

\textbf{What happened}: IBM's Watson Health represented the largest documented corporate investment in medical AI---and its most public failure. The MD Anderson Cancer Center partnership consumed \$62 million over four years before termination in 2017 with zero clinical deployment \citep{nytimes2021watson, forbes2017watson}. Watson for Genomics was discontinued in 2020; Watson for Oncology was shelved in 2021. IBM ultimately sold the Watson Health division, representing estimated total investment approaching \$4 billion.

\textbf{Failure mode}: Tacit knowledge gap. Oncology requires synthesis across patient history, imaging, genomics, clinical presentation, and treatment response---sphere architecture that pattern extraction could not replicate despite massive investment.

\textbf{Thermodynamic interpretation}: Dr. Norman Sharpless, former head of UNC's cancer center: ``We thought it would be easy, but it turned out to be really, really hard'' \citep{nytimes2021watson}. When Memorial Sloan Kettering tested Watson for Oncology, the system produced recommendations oncologists had already identified, demonstrating no added value beyond what sphere-trained physicians achieved through direct examination \citep{stat2017watson}.

\textbf{Ethical stake}: Opportunity cost of misdirected healthcare AI investment; patient expectations raised by marketing that capabilities could not meet.

\subsubsection{Case: McDonald's Drive-Thru AI---Acoustic Complexity}

\textbf{What happened}: In June 2024, McDonald's terminated its three-year partnership with IBM on Automated Order Taking (AOT) technology, removing the system from over 100 U.S. restaurants by July 26, 2024 \citep{restaurantbusiness2024}. Viral documentation showed the system ordering 260 Chicken McNuggets when customers requested a single meal, merging adjacent drive-thru lane orders, and continuing to add items while customers pleaded ``Stop! Stop! Stop!''

\textbf{Failure mode}: Collective tacit knowledge gap. Drive-thru ordering requires real-time adaptation to ambient noise, speech variation, order modifications, and social context---\citet{collins2010}'s collective tacit knowledge that emerges from participation in human communicative contexts.

\textbf{Thermodynamic interpretation}: CNBC reported the technology ``had issues interpreting different accents and dialects, which affected order accuracy'' \citep{cnbc2024mcdonalds}. Pattern extraction from clean training audio could not replicate the energetic investment of lived linguistic experience. McDonald's Chief Restaurant Officer Mason Smoot acknowledged in the franchisee memo the need to ``explore voice ordering solutions more broadly'' \citep{restaurantbusiness2024}.

\textbf{Ethical stake}: Customer frustration; reputational damage; employee displacement for systems that could not perform advertised functions.

\subsubsection{Case: Legal AI Hallucination---Confident Falsity as Architecture}

\textbf{What happened}: On June 22, 2023, Judge P. Kevin Castel imposed \$5,000 sanctions on attorneys Steven Schwartz and Peter LoDuca for submitting a legal brief containing six fabricated cases generated by ChatGPT \citep{fortune2023chatgpt}. The invented precedents---including \textit{Varghese v. China Southern Airlines}---contained bogus quotes and fictional judicial opinions. When questioned, ChatGPT assured Schwartz the cases were real and available on legal databases.

\textbf{Failure mode}: Hallucination as architectural feature. Large language models generate statistically plausible sequences without access to truth values; legal reasoning requires synthesis across precedent, statutory interpretation, and factual application.

\textbf{Thermodynamic interpretation}: Legal analyst Damien Charlotin documented over 660 cases of AI-generated false citations by December 2025, with incidents accelerating to 4--5 new cases daily. Stanford HAI research found legal LLMs hallucinate at rates varying from 69\% to 88\% for specific legal research tasks \citep{stanfordhai2024}. Judge Castel's ruling: ``Technological advances are commonplace\ldots But existing rules impose a gatekeeping role on attorneys to ensure the accuracy of their filings.''

\textbf{Ethical stake}: Professional responsibility; access to justice compromised by unreliable legal research; court system integrity; sanctions regime emerging to address systemic risk.

\subsection{Governance Failures: Extraction as Institutional Strategy}

The following cases illustrate how automation enables institutional strategies that would be ethically or legally impermissible if executed through transparent human decision-making. These represent governance failures that persist---and may even be exacerbated---regardless of underlying model capability.

\subsubsection{Case: UnitedHealth nH Predict---Algorithmic Denial Optimization}

\textbf{What happened}: A federal lawsuit filed in November 2023 alleges UnitedHealth deployed the nH Predict algorithm with a 90\% error rate---nine of ten appealed denials ultimately reversed---to systematically deny elderly patients coverage for post-acute care \citep{cbsnews2023united}. STAT News investigation documented UnitedHealthcare's denial rate for post-acute care rising from 10.9\% in 2020 to 22.7\% in 2022, coinciding with nH Predict deployment \citep{statnews2023medicare}.

\textbf{Failure mode}: Incentive misalignment. The algorithm was not designed to replicate physician judgment but to override it in service of cost reduction. The lawsuit alleges UnitedHealth set goals to keep skilled nursing facility stays within 1\% of algorithm predictions, with employees ``disciplined and terminated'' for deviation \citep{cbsnews2023united}.

\textbf{Thermodynamic interpretation}: The extraction calculus exploits asymmetry: ``Despite the high error rate, Defendants continue to systemically deny claims\ldots because they know that only a tiny minority of policyholders (roughly 0.2\%) will appeal denied claims.'' The system extracts savings from patients lacking resources to contest. A February 2025 court ruling allowed the lawsuit to proceed, with Judge John Tunheim describing the claims process as ``futile'' with likelihood of causing ``irreparable injury'' \citep{courthousenews2025united}.

\textbf{Ethical stake}: Patient welfare; informed consent to coverage limitations; due process in benefits determination; exploitation of vulnerable populations' limited contestation capacity.

\subsubsection{Case: COMPAS and Predictive Policing---Bias Feedback Loops}

\textbf{What happened}: The COMPAS recidivism algorithm, used in sentencing across multiple states, was found by ProPublica to ``falsely label [Black defendants] as future criminals at twice the rate as white defendants'' \citep{propublica2016compas}. In \textit{State v. Loomis} (2016), the Wisconsin Supreme Court allowed continued use but mandated warnings about proprietary opacity, lack of state-specific validation, and racial disproportionality \citep{wisconsin2016loomis}.

\textbf{Failure mode}: Bias feedback loop. Algorithms trained on historically biased arrest data extract that bias into predictions, which inform policing decisions, which generate arrest data reinforcing original bias. Even without explicit race variables, socioeconomic proxies replicate protected-class discrimination.

\textbf{Thermodynamic interpretation}: MIT Technology Review documented how mass arrest events produced data that disproportionately targeted communities in future predictions \citep{mittechreview2020}. Los Angeles activist Hamid Khan sought audit of PredPol; the inspector general declared the task ``impossible because the tool was so complicated.'' U.S. Senators wrote to DOJ: ``mounting evidence indicates that predictive policing technologies do not reduce crime\ldots Instead, they worsen the unequal treatment of Americans of color by law enforcement'' \citep{senate2024predictive}.

\textbf{Ethical stake}: Due process; equal protection; opacity preventing meaningful contestation; state power amplified through unaccountable algorithmic systems.

\subsubsection{Case: Apple Card---Proxy Discrimination}

\textbf{What happened}: Apple Card's credit algorithm granted Danish entrepreneur David Heinemeier Hansson a credit limit 20 times higher than his wife despite her superior credit score; Apple co-founder Steve Wozniak confirmed similar disparities \citep{bbc2019apple}. The Consumer Financial Protection Bureau imposed over \$89 million in combined penalties and redress in October 2024 \citep{cfpb2024apple}.

\textbf{Failure mode}: Proxy discrimination. The algorithm produced discriminatory outcomes without explicit use of gender as input variable---socioeconomic and behavioral proxies replicated protected-class bias.

\textbf{Thermodynamic interpretation}: New York DFS Superintendent Linda Lacewell: ``Any algorithm that intentionally or not results in discriminatory treatment of women or any other protected class violates New York law'' \citep{bbc2019apple}. The extraction of creditworthiness signals from historical data reproduced historical discrimination patterns.

\textbf{Ethical stake}: Fair lending; gender equality; accountability for algorithmic outcomes regardless of intent; regulatory enforcement of anti-discrimination law in automated systems.

\subsubsection{Case: Air Canada Chatbot---Liability for Extraction Systems}

\textbf{What happened}: On February 14, 2024, the British Columbia Civil Resolution Tribunal found Air Canada liable for negligent misrepresentation after its chatbot incorrectly told Jake Moffatt he could apply for bereavement fares retroactively \citep{bccrt2024moffatt}. Air Canada argued it ``cannot be held liable for information provided by one of its agents, servants, or representatives---including a chatbot,'' suggesting the chatbot was ``a separate legal entity that is responsible for its own actions.''

\textbf{Failure mode}: Accountability evasion attempt. The airline sought to externalize liability to its own extraction system.

\textbf{Thermodynamic interpretation}: Tribunal member Christopher C. Rivers found this argument ``remarkable'': ``It should be obvious to Air Canada that it is responsible for all the information on its website. It makes no difference whether the information comes from a static page or a chatbot.'' The ruling establishes that organizations cannot outsource responsibility to extraction systems; the energy investment required for accuracy remains with the deploying entity.

\textbf{Ethical stake}: Consumer protection; corporate accountability; precedent preventing ``the algorithm did it'' defense; representation accuracy in automated customer interactions.

\subsection{The Reversal Pattern: Documented Course Corrections}

\subsubsection{Case: Klarna---``Too Far, Too Fast''}

\textbf{What happened}: In February 2024, Klarna announced its AI assistant handled 2.3 million conversations in the first month, performing ``the equivalent work of 700 full-time agents'' with claimed \$40 million annual savings \citep{klarna2024ai}. By May 2025, CEO Sebastian Siemiatkowski reversed course, acknowledging the company had gone ``too far, too fast'' as customer satisfaction declined \citep{bloomberg2025klarna}.

\textbf{Failure mode}: Premature extraction. Cost optimization degraded the judgment, empathy, and contextual adaptation that complex customer service requires.

\textbf{Thermodynamic interpretation}: Siemiatkowski told Bloomberg: ``From a brand perspective, a company perspective\ldots I just think it's so critical that you are clear to your customer that there will be always a human if you want.'' Klarna launched a pilot program rehiring human agents. Julie Geller of Info-Tech Research Group: ``AI should augment human agents---not replace them. Automate the routine\ldots but always ensure customers have a clear, easy path to a human'' \citep{cxdive2025klarna}.

\textbf{Ethical stake}: Customer experience; employee displacement and subsequent rehiring; organizational learning about automation limits.

\subsubsection{Case: Commonwealth Bank of Australia---Two-Month Collapse Cycle}

\textbf{What happened}: In June 2025, Commonwealth Bank introduced a chatbot to handle customer queries, laying off 45 workers. The system diverted 2,000 calls weekly. By August 2025---within two months---the bank reversed course and began rehiring \citep{techtarget2025cba}.

\textbf{Failure mode}: Complexity underestimation. The chatbot could not handle complex queries representing a larger workload portion than anticipated.

\textbf{Thermodynamic interpretation}: Extraction of ``simple'' queries left concentrated complexity that the chatbot could not navigate. The two-month reversal demonstrates how quickly brittleness manifests when extraction systems encounter genuine cognitive complexity.

\textbf{Ethical stake}: Employee welfare; customer service quality; organizational decision-making about automation deployment.

\subsection{The Micro-Credential Delusion}

\citet{lumina2025} reports that 96\% of employers say micro-credentials strengthen job applications, with 90\% offering 10--15\% higher starting salaries. Yet this represents preference for granular verification over diluted degrees, not evidence of capability development.

\citet{ha2022}'s systematic review reveals the measurement gap: while 13 of 14 studies show positive outcomes, these measure satisfaction and employment signals rather than demonstrated competence. \citet{joshi2019} provides calibration: coding bootcamps help non-technical graduates enter technology fields but provide ``minimal benefit for those already holding technical degrees.''

The credentialing trajectory in the cases reviewed here:

\begin{itemize}
\item Google Career Certificate: 3--6 months, single skill vector
\item Coursera Specialization: 4--8 weeks, narrow domain
\item LinkedIn Learning Path: 5--10 hours, procedural fragments
\item Energy investment per credential: approaching thermodynamic minimum
\item Cumulative synthesis capacity developed: not measured, likely negligible
\end{itemize}

Micro-credentials optimize for signaling efficiency, not cognitive development---the educational analog of extraction without maintenance investment.

\subsection{Pattern Synthesis: What the Evidence Suggests}

The cases reviewed here, while illustrative rather than representative, reveal consistent patterns suggestive of broader dynamics warranting systematic study:

\textbf{Capability failures} share common characteristics: performance measured on training distributions does not transfer to deployment complexity; tacit knowledge gaps prove resistant to extraction regardless of investment scale; maintenance entropy accumulates in systems optimized for initial deployment rather than sustained operation.

\textbf{Governance failures} exploit automation for institutional strategies that would face resistance if executed transparently: denial optimization, bias amplification, accountability evasion. These persist---and may intensify---regardless of underlying model capability.

Both failure modes share a thermodynamic signature: value extraction from cognitive systems without corresponding investment in maintenance, validation, oversight, or human capability preservation. The 42\% abandonment rate, the 80\%+ project failure rate, the documented reversals---these suggest not technological immaturity to be solved through iteration, but structural constraints on what extraction-based approaches can achieve.

The evidence does not support claims that AI ``cannot work'' or that automation is inherently problematic. It suggests, rather, that deployment strategies treating human cognitive capability as a resource to be extracted---rather than a system requiring ongoing investment---encounter predictable brittleness when complexity exceeds training distribution boundaries or when accountability structures fail to constrain institutional incentives.

The thermodynamic framing offers analytic purchase on these patterns: systems that extract without investing in maintenance face entropy accumulation; capabilities not actively sustained through energy investment decay; complexity requires cognitive architecture that educational and organizational vectorization systematically depletes.

This framing does not explain everything. Institutional incentives, regulatory capture, labor market dynamics, and political economy all shape AI deployment outcomes through mechanisms not reducible to thermodynamic analogy. But the energy-budget constraint---the principle that cognitive capabilities require sustained investment and cannot be extracted without entropic consequence---provides a unifying lens for understanding why so many high-profile deployments fail in structurally similar ways.

The physics, as analytic constraint rather than metaphysical claim, does not negotiate. The evidence, in the cases reviewed here, accumulates. The pattern, suggestive if not definitive, warrants serious attention.
